\begin{abstract}

Segmentation models in medical imaging produce predictions for every pixel, but
clinical deployment requires knowing \emph{which} predictions to trust. Current
practice treats uncertainty estimation as an afterthought: a model segments an
image, an uncertainty map is computed, and then both are handed to a clinician
with no systematic policy for when to accept, flag, or defer. We formalize this
gap as a selective prediction problem and study how different uncertainty
sources and deferral strategies interact when the goal is not better
segmentation per se, but better \emph{decisions about segmentation}.

Working on retinal vessel analysis (DRIVE, STARE, CHASE\_DB1), we compare two
uncertainty estimation methods, Monte Carlo (MC) Dropout and Test-Time
Augmentation (TTA), under three increasingly sophisticated deferral policies:
global thresholding, image-adaptive thresholding, and a confidence-aware
strategy that jointly weighs uncertainty and prediction confidence. We also
evaluate temperature scaling as a calibration step and find, contrary to
common assumption, that it does not improve downstream decision quality.

Our central result: TTA paired with adaptive deferral reduces segmentation
error by approximately 80\% while deferring only 25\% of pixels. This
demonstrates that the right combination of uncertainty source and deferral
policy can eliminate the vast majority of segmentation errors at modest cost.
In the unified comparison bundle, TTA reaches an uncertainty AUROC of 0.881,
against 0.722 for MC Dropout, while running about 3.1$\times$ faster
(0.86\,s vs.\ 2.65\,s per image). Confidence-aware deferral offers the best
low-budget efficiency: it removes about 55\% of the error while deferring only
12\% of pixels. Temperature scaling does not improve the downstream operating
points in this repository and worsens the tracked ECE values for both methods
in the main comparison bundle. Cross-dataset evaluation of the MC Dropout
pipeline shows that uncertainty quality transfers reasonably well under domain
shift: uncertainty AUROC remains above 0.83 on STARE and CHASE\_DB1 while Dice
falls from 0.791 on DRIVE to about 0.726.

The practical conclusion: uncertainty methods should be chosen and evaluated
by what they enable downstream (error reduction, deferral efficiency), not by
calibration metrics. The best configuration in our experiments, TTA with
confidence-aware deferral, requires no model retraining, no auxiliary
networks, and keeps the review budget low.

\end{abstract}
